\documentclass[11pt]{article}
\usepackage[margin=1in]{geometry}
\usepackage{amsmath,amssymb,graphicx,booktabs,hyperref,xcolor}
\title{Computational Replication Report:\\ \'Smejkal, Sinova \& Jungwirth (2022),\\ \emph{Emerging research landscape of altermagnetism} (arXiv:2204.10844)}
\author{Ollie (OpenClaw @ Argonne) --- TEXTURES-100 Wave 4}
\date{2026-07-19}
\begin{document}
\maketitle

\section{Paper summary}
This Perspective delimits \emph{altermagnetism} as a third basic collinear magnetic
phase distinct from ferromagnetism and antiferromagnetism. Altermagnets combine a
\emph{compensated} (zero net moment, AFM-like) collinear spin order with
\emph{nonrelativistic spin-split} electronic bands (FM-like), reconciled by an
unconventional spin-group symmetry formalism. The two spin sublattices are connected
by a proper/improper crystal \emph{rotation} rather than translation or inversion,
producing an even-parity ($d$-, $g$-, $i$-wave) momentum-dependent spin splitting with
nodal lines and sign-alternating lobes, so that the Brillouin-zone integral of the spin
splitting vanishes while the local splitting is large.

\section{Claims addressed}
\begin{table}[h]\centering\small
\begin{tabular}{clccc}
\toprule
ID & Claim & Type & Testable? & Tested? \\
\midrule
C1 & Compensated ($M_{\rm net}=0$) yet nonrelativistically spin-split bands & conceptual/model & yes & yes \\
C2 & Symmetry rule: sublattices connected by rotation, not translation/inversion & symmetry & yes & yes \\
C3 & Even-parity $d$-wave splitting: nodal lines, 4 alternating lobes, BZ-compensated & model & yes & yes \\
C4 & Material zoo (RuO$_2$, MnTe, CrSb\ldots) from ab initio & DFT & yes & \textcolor{red}{no (out of scope)} \\
\bottomrule
\end{tabular}
\end{table}

\section{Method}
We use a minimal single-orbital square-lattice tight-binding altermagnet
(reusing the Wave-3 \texttt{sasioglu2026} $d$-wave template). The spin-$\sigma$ dispersion is
\[
E_\sigma(\mathbf k) = -2t(\cos k_x+\cos k_y) - 2\sigma\, t_{\rm AM}(\cos k_x-\cos k_y),
\]
giving the $d$-wave spin splitting
$\Delta(\mathbf k)=E_\uparrow-E_\downarrow=-4t_{\rm AM}(\cos k_x-\cos k_y)$,
with $t=1$, $t_{\rm AM}=0.30$. Tools: Python 3.13, NumPy, Matplotlib
(\texttt{code/smejkal2022\_replication.py}). Diagnostics computed on a $401\times401$
$\mathbf k$-grid: BZ-averaged splitting, half-filling net moment $n_\uparrow-n_\downarrow$,
the $C_4$+spin-flip symmetry residual $\max|E_\uparrow(\mathbf k)-E_\downarrow(C_4\mathbf k)|$,
the translation residual, the diagonal nodal amplitude, and the number of sign changes of
$\Delta$ around a $\Gamma$-centered circle.

\section{Results vs.\ paper}
\begin{table}[h]\centering\small
\begin{tabular}{lccc}
\toprule
Quantity & Expected (paper concept) & This work & Verdict \\
\midrule
BZ-avg spin splitting & $0$ (compensated) & $0.00$ & \checkmark \\
Net moment (half filling) & $0$ & $0.00$ & \checkmark \\
Local max $|\Delta|$ & large (FM-like) & $2.40\,t$ & \checkmark \\
$C_4$+spin-flip residual & $0$ (rotation connects sublattices) & $0.0$ & \checkmark \\
Translation residual & large (NOT AFM-translation) & $2.40$ & \checkmark \\
Diagonal node $|\Delta|$ & $0$ (nodal lines) & $0.00$ & \checkmark \\
Sign changes around $\Gamma$ & $4$ ($d$-wave) & $4$ & \checkmark \\
\bottomrule
\end{tabular}
\end{table}

\section{Per-claim what worked / what didn't}
\textbf{C1 (worked).} $M_{\rm net}=0$ at half filling coexisting with $\max|\Delta|=2.4t$
reproduces the FM--AFM dichotomy exactly. \textbf{C2 (worked).} The $C_4$+spin-flip residual
is machine-zero while the translation residual is $2.4$, numerically demonstrating the paper's
identification rule. \textbf{C3 (worked).} Nodes on $k_x=\pm k_y$ and exactly 4 sign changes
confirm $d$-wave. \textbf{C4 (not attempted).} Material-specific ab initio band structures
require DFT and are out of scope for a CPU-only model replication.

\section{Critique of the replication}
This is a \emph{conceptual/model} replication of a review, not a numerical match to a single
figure. Strength: the symmetry and node structure that anchor the entire Perspective are
reproduced to machine precision, and the LLM-judge scored REPLICATED (coverage 9/10,
agreement 9/10). Weaknesses honestly noted: (i) the model is imposed, not derived --- $t_{\rm AM}$
is an input, so we do not prove altermagnetism is the energetic ground state; (ii) only the
$d$-wave case is shown, whereas the paper spans $g$- and $i$-wave materials; (iii) all energies
are in units of $t$, so there is no absolute-meV comparison to any compound; (iv) SOC, magnons,
and transport (Secs III.C--D) are outside this minimal electronic model. None of these undercut
the central claims tested, but they bound the scope.

\section{Verdict}
\textbf{REPLICATED} (representative model + symmetry classification). C1--C3 reproduced to
machine precision; C4 (material DFT) out of scope and documented as method-limited, not a failure.

\section*{Open Questions}
See \texttt{report/open\_questions.json}. Q1: multi-orbital / higher-wave (g,i) material fidelity.
Q2: SOC-lifted nodes and anomalous Hall magnitude. Q3: self-consistent interacting origin
(Hubbard mean field). Q4: chiral magnon splitting in the same model. Q5: cleanest transport
signature distinguishing altermagnet from AFM.

\end{document}
